\documentclass[12pt,a4paper]{article}
% 12pt:     Tamaño de fuente
% a4paper: Tamaño de papel
% article:  Formato artículo, está dividido por secciones y subsecciones. Otros formatos: report, letter, book, etc.

\usepackage[spanish, es-tabla]{babel}
% Configura en español los nombres de figuras, tablas, etc.
\usepackage[utf8]{inputenc}                 % Permite el uso directo de tildes y otros.
\usepackage{amsmath,amssymb}
% Añade símbolos matemáticos para las ecuaciones.
\usepackage{parskip}
% Elimina la sangría
\usepackage{geometry}
% Permite modificar los márgenes de la página a gusto.
\usepackage{multicol}
% Ordena el cuerpo del documento en 2 o más columnas.
\usepackage{graphicx}
% Para poner imágenes.
\usepackage{enumitem}
% Para modificar los tags de los enumerate.
\usepackage{array}
% Para ordenar texto y ecuaciones.
\usepackage{booktabs}
% Proporciona las reglas \toprule, \midrule y \bottomrule para tablas.
\usepackage{url}
% Permite usar url que redirigen a la página web.
\usepackage{hyperref}
% Permite que se puedan clickear las urls, citas y referencias.
\usepackage{float}

\spanishdecimal{.} % punto decimal
\geometry{margin=2cm} % margenes de la hoja
\renewcommand{\baselinestretch}{1.3} % tamaño de las filas de las tablas

\newcommand{\uveci}{\hat{\imath}} % vector unitario i
\newcommand{\uvecj}{\hat{\jmath}} % vector unitario j
\newcommand{\uveck}{\hat{k}} % vector unitario k

\begin{document}

\begin{center}
{\Large\textbf{INFORME N°3}}

{\LARGE\textbf{Ley de Hooke y constante elástica de un resorte}}

\rule{\textwidth}{0.5pt}

\begin{tabular}{ll}
    \textbf{Estudiantes:} & Nombre Apellido \\
                          & Nombre Apellido \\
                          & Nombre Apellido \\
                          & Nombre Apellido \\[0.2cm]
    \textbf{Profesora:}   & Nombre Apellido \\
    \textbf{Ayudantes:}    & Nombre Apellido \\
                           & Nombre Apellido\\[0.2cm]
    \textbf{Fecha de entrega:} & DD de mes de AAAA
\end{tabular}

\vspace{0.4cm}
\rule{\textwidth}{0.5pt}

\end{center}

\section{Introducción}

El estudio se centra en el comportamiento elástico de los cuerpos deformables, específicamente en la respuesta de un resorte ante la acción de una fuerza externa. Cuando un resorte es estirado o comprimido, experimenta una fuerza recuperadora que tiende a devolverlo a su longitud natural. Experimentalmente, la importancia de analizar la relación entre la fuerza aplicada ($F$) y la elongación ($x$) radica en comprobar la Ley de Hooke dentro del límite elástico, donde ambas magnitudes guardan una relación de proporcionalidad directa ($F = kx$). Esto permite verificar de forma empírica si el comportamiento mecánico del resorte se ajusta a los modelos teóricos ideales mediante el análisis gráfico y el uso de ajustes lineales.
La constante elástica (k) representa una propiedad característica de cada resorte que cuantifica su oposición a ser deformado. Experimentalmente, su determinación a través de la pendiente de un ajuste lineal entre la fuerza y la elongación, es fundamental porque permite caracterizar y diferenciar la rigidez entre distintos resortes; un valor de k mayor indica un resorte más rígido y difícil de estirar, mientras que uno menor corresponde a uno más flexible.
En muchas aplicaciones físicas e ingenieras, los sistemas mecánicos están compuestos por múltiples elementos elásticos combinados. La posibilidad de caracterizar un sistema de resortes conectados en serie mediante una constante elástica equivalente ($k_{1+2}$) es de gran importancia experimental y teórica, ya que simplifica el análisis del conjunto. Permite predecir y contrastar cómo se comportaría el sistema completo bajo una carga global aplicando las leyes de asociación (como la relación inversa para resortes en serie), evaluando si el modelo teórico representa con precisión el comportamiento colectivo de los componentes reales.

\section{Fundamento teórico}

Deduzcan y presenten las ecuaciones de:

\begin{itemize}
    \item Relación entre fuerza elástica y elongación (ley de Hooke)
    \item Condición de equilibrio que lleva a determinar la fuerza
\end{itemize}

Indique además, sin necesidad de deducirlas completas, las ecuaciones de:

\begin{itemize}
    \item Constante elástica equivalente del sistema de resortes conectados en serie
    \item Cálculo de propagación de error
\end{itemize}

En esta sección debe presentar las relaciones físicas y las deducciones matemáticas necesarias para analizar los resultados. No basta con enumerar fórmulas, sino que se debe explicar de dónde provienen y qué representan sus variables.

\section{Procedimiento experimental}

\subsection{Montaje} \label{Montaje}

Describan el montaje experimental, con las características de cada uno de sus materiales y cómo están conectados entre sí.

Identifiquen, como lista o como párrafo, los instrumentos utilizados para medir cada magnitud y las resoluciones correspondientes.

\subsection{Parte 1 - Determinar constante elástica de cada resorte} \label{Procedimiento-Parte-1}

Describan el procedimiento general utilizado para la realización de la medición de las longitudes naturales, las fuerzas aplicadas y las longitudes deformadas de ambos resortes por separado. Describan también el procedimiento utilizado para recopilar los datos y calcular su promedio, desviación estándar y error absoluto. Pueden incluir observaciones cuantitativas y cualitativas.

\subsection{Parte 2 -Determinar constante elástica de resortes en serie} \label{Procedimiento-Parte-2}

Describan el procedimiento general utilizado para la realización de la medición de las longitudes naturales, las fuerzas aplicadas y las longitudes deformadas de ambos resortes conectados en serie. Describan también el procedimiento utilizado para recopilar los datos y calcular su promedio, desviación estándar y error absoluto. Pueden incluir observaciones cuantitativas y cualitativas.


\section{Resultados}

\subsection{Constantes elásticas de los resortes individuales}

\begin{table}[H]
    \centering
    \caption{Tabla configuraciones resorte 1}
    \label{tab:tabla1}
    \begin{tabular}{c c c c c c}
        \toprule
        Nº & $\bar{L}_0$ & $\bar{x}$ & $\Delta{x}$ & $F$ & $\Delta{F}$ \\
        \midrule
        1 & $0,144 m$ & $0,089 m$ & $\pm{0,001} m$ & $2,45 N$ & $\pm{0,001} N$ \\
        2 & $0,144 m$ & $0,173 m$ & $\pm{0,001} m$ & $4,90 N$ & $\pm{0,001} N$ \\
        3 & $0,144 m$ & $0,260 m$ & $\pm{0,001} m$ & $7,36 N$ & $\pm{0,001} N$ \\
        4 & $0,144 m$ & $0,344 m$ & $\pm{0,001} m$ & $9,83 N$ & $\pm{0,001} N$ \\
        5 & $0,144 m$ & $0,428 m$ & $\pm{0,001} m$ & $12,3 N$ & $\pm{0,001} N$ \\
        6 & $0,144 m$ & $0,515 m$ & $\pm{0,001} m$ & $14,7 N$ & $\pm{0,001} N$ \\
         \bottomrule
    \end {tabular}
\end{table}

\begin{table}[H]
    \centering
    \caption{Tabla configuraciones resorte 2}
    \label{tab:tabla2}
    \begin{tabular}{c c c c c c}
        \toprule
        Nº & $\bar{L}_0$ & $\bar{x}$ & $\Delta{x}$ & $F$ & $\Delta{F}$ \\
        \midrule
        1 & $0,194 m$ & $0,093 m$ & $\pm{0,001} m$ & $2,45 N$ & $\pm{0,001} N$ \\
        2 & $0,194 m$ & $0,188 m$ & $\pm{0,001} m$ & $4,90 N$ & $\pm{0,001} N$ \\
        3 & $0,194 m$ & $0,278 m$ & $\pm{0,001} m$ & $7,36 N$ & $\pm{0,001} N$ \\
        4 & $0,194 m$ & $0,368 m$ & $\pm{0,001} m$ & $9,83 N$ & $\pm{0,001} N$ \\
        5 & $0,194 m$ & $0,459 m$ & $\pm{0,001} m$ & $12,3 N$ & $\pm{0,001} N$ \\
        6 & $0,194 m$ & $0,547 m$ & $\pm{0,001} m$ & $14,7 N$ & $\pm{0,001} N$ \\
        \bottomrule
    \end{tabular}
\end{table}

\begin{figure}[H]
    \centering
    \includegraphics[width=1\textwidth]{resorte1_grafico.png}
    \caption{gráfico F vs. x del resorte 1.}
    \label{fig:grafico1}
\end{figure}

\begin{figure}[H]
    \centering
    \includegraphics[width=0.8\textwidth]{resorte2_grafico.png}
    \caption{gráfico F vs. x del resorte 2.}
    \label{fig:grafico2}
\end{figure}

Para el resorte 1, la pendiente de la recta asignada es de $a=29,46$, el intercepto con el eje F es de $b=-0,39$ y la constante k del resorte 1 es de $k=29,46 N/m$.
Para el resorte 2, la pendiente de la recta asignada es de $a=27,65$, el intercepto con el eje F es de $b=-0,42$ y la constante k del resorte 2 es de $k=27,65 N/m$.
En ambos casos, tener presente que las barras de error son bastante poco visibles debido a que es del orden de los milimetros y el gráfico está en el orden de metros.

\subsection{Constante elástica equivalente del sistema en serie}

\begin{table}[H]
    \centering
    \caption{Tabla configuraciones resortes en serie}
    \label{tab:tabla3}
    \begin{tabular}{c c c c c c}
        \toprule
        Nº & $\bar{L}_0$ & $\bar{x}$ & $\Delta{x}$ & $F$ & $\Delta{F}$ \\
        \midrule
        1 & $0,388 m$ & $0,177 m$ & $\pm{0,001} m$ & $2,45 N$ & $\pm{0,001} N$ \\
        2 & $0,388 m$ & $0,355 m$ & $\pm{0,001} m$ & $4,90 N$ & $\pm{0,001} N$ \\
        3 & $0,388 m$ & $0,530 m$ & $\pm{0,001} m$ & $7,36 N$ & $\pm{0,001} N$ \\
        4 & $0,388 m$ & $0,705 m$ & $\pm{0,001} m$ & $9,83 N$ & $\pm{0,001} N$ \\
        5 & $0,388 m$ & $0,880 m$ & $\pm{0,001} m$ & $12,3 N$ & $\pm{0,001} N$ \\
        6 & $0,388 m$ & $1,05 m$ & $\pm{0,001} m$ & $14,7 N$ & $\pm{0,001} N$ \\
        \bottomrule
    \end{tabular}
\end{table}

\begin{figure}[H]
    \centering
    \includegraphics[width=1\textwidth]{resortes_en_serie_grafico.png}
    \caption{gráfico F vs. x de los resortes en serie.}
    \label{fig:graficoserie}
\end{figure}

Para los resortes en serie, la pendiente de la recta asignada es de $a=14,35$, el intercepto con el eje F es de $b=-0,35$ y la constante k de los resortes en serie es igual a $k=14,35 N/m$.
Tener en cuenta que las barras de error son muy poco visibles debido a que el error está en el orden de los milímetros y el gráfico está en el orden de los metros.

\subsection{Comparación entre resultados experimentales y predicción teórica}

    \begin{table}[H]
    \centering
        \caption{tabla resumen}
        \label{tab:tabla4}
        \begin{tabular}{c c c}
        \toprule
            Sistema & Pendiente a & Intercepto b\\
        \midrule
            Resorte 1 & 29,46 & -0,39 \\
            Resorte 2 & 27,65 & -0,42 \\
            R1 + R2 (serie) & 14,35 & -0,35 \\
        \bottomrule
        \end{tabular}
    \end{table}

Para encontrar nuestra predicción teórica de la constante de $k_{1+2 teo}$ usaremos la equivalencia $\frac{1}{k_{1+2 teo}} = \frac{1}{k_1} + \frac{1}{k_2}$. 

$\frac{1}{k_{1+2 teo}} = \frac{1}{29,46} + \frac{1}{27,65} = 14,26$N/m.

Por ende, si lo comparamos con $k_{1+2 exp} = 14,35$, N/m tenemos que la constante teorica es bastante acertado con respecto a la constante experimental.


\section{Análisis y conclusiones}

En esta sección, resuman brevemente lo realizado en el laboratorio. Indiquen qué se hizo y cuál fue el propósito de cada actividad, haciendo referencia a las secciones, tablas y ecuaciones correspondientes cuando sea necesario.

Luego del resumen de las actividades, respondan las preguntas de análisis planteadas en las tres partes de la guía. Las respuestas a las preguntas de análisis deben estar integradas en el texto, conectando ideas mediante conectores lógicos.

Finalmente, incluyan las principales conclusiones a las que llegaron como grupo. Cuando corresponda, indiquen también las principales limitaciones o fuentes de incertidumbre que hayan identificado durante el trabajo.

Las conclusiones deben basarse en los resultados obtenidos por el grupo. \textbf{No incluyan citas bibliográficas en esta sección}.

\textbf{\underline{Extensión máxima del informe:} 15 páginas}.

\begin{thebibliography}{99}

\bibitem{libro} Nombre del autor. \textit{Título del libro}. Año.

\bibitem{apunte} Nombre del profesor. \textit{Nombre del curso}. Semestre y año en el que fue entregado el apunte.

\bibitem{web} TikTok de @usuario. Disponible en: \url{https://www.tiktok.com/}.

\end{thebibliography}

\end{document}